% !TEX root = tarasov-qasp-zh.tex
% 本章首尾均无有效 sec 块：第 IV 部分（静电学）的 sec-i 扉图与格言在英文源 ch-22～ch-25 尾部
% 均为 % 注释的死代码副本（含本文件尾部），未搬入。

\chapter[静电场是如何描述的？]{静电场是如何描述的？\label{ch-23}}

\begin{dialogue}

\paragraph{}
\tchr 这样，我们就来继续上一节以如下问题开始的讨论：\hlred{“静电场是如何描述的？”}

\stdb 静电场（electrostatic field）用一种矢量的力的特征量来描述，它叫做电场强度（intensity of the electric field）。在场中每一点，电场强度 $E$ 都有确定的方向和数值。如果我们在场中从一点移动到另一点，且移动的方向始终与电场强度矢量的方向相切，这样得到的路径就称为场的电场线（lines of force）。用电场线来直观地描绘一个场是非常方便的。
\pnote{电场强度（简称场强）是人教版必修第三册第九章“电场 电场强度”一节的核心概念，定义为 $E=F/q$，方向规定为正电荷受力的方向——正是式~\eqref{eq-129} 的内容。lines of force 直译为“力线”，现行教材通称电场线：曲线上每一点的切线方向就代表该点场强的方向，原书“让移动方向始终与场强矢量相切所得的路径”是同一定义的另一种说法。}

\tchr 好。现在我们把问题谈得更具体些。相距为 $r$ 的两个电荷 $q_{1}$ 与 $q_{2}$ 之间的库仑力（Coulomb force）可以写成
\begin{equation}%
F_{e} = \frac{q_{1} q_{2}}{r^{2}}
\label{eq-127}
\end{equation}
的形式。它可以改写为
\begin{equation}%
E(r) =  \dfrac{q_{1}}{r^{2}}
\label{eq-128}
\end{equation}
\begin{equation}%
F_{e} = E(r) q_{2}
\label{eq-129}
\end{equation}
式~\eqref{eq-128} 表明，电荷 $q_{1}$ 在自己周围建立场，在离电荷 $r$ 处，这个场的强度等于 $\dfrac{q_{1}}{r^{2}}$。式~\eqref{eq-129} 表明，这个场对位于离电荷 $q_{1}$ 为 $r$ 处的电荷 $q_{2}$ 作用一个力 $E(r) q_{2}$。之所以能把式~\eqref{eq-127} 写成这种形式，是因为引入了一个“中间人”——量 $E$，即场的特征量。请试着确定式~\eqref{eq-127}、\eqref{eq-128} 和~\eqref{eq-129} 的适用范围。
\pnote{本节公式不含 $1/(4\pi\varepsilon_{0})$ 因子：原书静电学部分沿用苏联教材惯用的高斯制（绝对静电单位制），真空中的库仑定律直接写作 $F=q_{1}q_{2}/r^{2}$，比例常数为 1，点电荷的场强相应为 $E=q/r^{2}$。人教版必修第三册第九章采用 SI：$F=kq_{1}q_{2}/r^{2}$，静电力常量 $k=1/(4\pi\varepsilon_{0})\approx 9.0\times 10^{9}\;\mathrm{N \cdot m^{2}/C^{2}}$，点电荷的场强为 $E=kQ/r^{2}$。两种单位制只差一个常数因子，场的叠加、电势差等一切方法性结论不变；但场强与电势的单位在两制中不同，数值不能直接互换。}

\stdb 式~\eqref{eq-127} 适用于两个点电荷（point charge）。这就是说，式~\eqref{eq-128} 和式~\eqref{eq-129} 的适用范围也相同，因为它们是从式~\eqref{eq-127} 得出来的。

\tchr 这只对式~\eqref{eq-127} 和式~\eqref{eq-128} 成立。式~\eqref{eq-129} 的适用范围要广泛得多。无论场 $E$ 是由什么建立的（一个点电荷、一组点电荷，还是任意形状的带电体），在所有情形下，这个场对电荷 $q_{0}$ 的作用力都等于该电荷与电荷 $q_{0}$ 所在处场强的乘积。式~\eqref{eq-129} 更普遍的形式具有如下的矢量形式
\begin{equation}%
\vec{F}_{e} = \vec{E}(r) \, q_{0}
\label{eq-130}
\end{equation}
其中箭头照例用来表示矢量。由式~\eqref{eq-130} 可以看出：若电荷 $q_{0}$ 为正，则在场的给定一点作用在电荷 $q_{0}$ 上的力，其方向与该点场强方向一致；若电荷 $q_{0}$ 为负，则力的方向与场强方向相反。

在这里可以体会到场这一概念自身的独立性。不同的带电体建立不同的静电场，但其中每一个场都按照同一规律~\eqref{eq-130} 作用于处在场中的电荷。要想求出作用在电荷上的力，必须先算出电荷所在点的场强。因此，会求一组电荷所建立的场的场强是很重要的。设有两个电荷 $q_{1}$ 和 $q_{2}$。对我们可能感兴趣的空间任一点，由其中每个电荷建立的场的场强，其大小和方向都不难求出。设在由矢量 $\vec{r}$ 给定的某一点，这些场强由矢量 $\vec{E}_{1}(\vec{r})$ 和 $\vec{E}_{2}(\vec{r})$ 描述。要求出点 $\vec{r}$ 处的合场强，必须把各个电荷单独产生的场强按矢量相加：
\begin{equation}%
\vec{E}(r) = \vec{E}_{1}(r) + \vec{E}_{2}(r)
\label{eq-131}
\end{equation}
我再说一遍：场强必须按矢量相加。（问学生甲）你听明白了吗？

\stda 明白，我知道场强是按矢量相加的。

\tchr 很好。那就来检验一下你在实践中运用这项知识的能力。\hlred{请画出两个等量异号电荷（$+q_{1}$ 与 $-q_{2}$）的场的电场线，并假定其中一个电荷（例如 $+q_{1}$）是另一个的几倍。}
\pnote{原句“两个等量异号电荷……假定其中一个是另一个的几倍”自相矛盾：既然等量，便无从“几倍”。按上下文，这里的任务当是画一对不等量异号点电荷的场的电场线——学生甲随即画出的等量异号情形（\figr{fig-86}）只是熟知的起点，教师随后才逐次改变两电荷的量值。“等量”疑为英文成书时沿用 equal and opposite 惯用语所致，待考原书。中文照译原句。}

\stda 恐怕我画不了。这样的场我们过去没有讨论过。

\tchr 那你们学过什么样的场呢？

\stda 我知道由两个等量点电荷建立的场的电场线图形是什么样子。这样的图我画在了\figr{fig-86} 里。

\begin{figure}%[-5cm]%
\centering
\includegraphics[width=0.8\textwidth]{figs/ch-23/fig-086a}
\caption[两个等量异号点电荷的场的电场线。]{两个等量异号点电荷的场的电场线。\label{fig-86}}
\end{figure}
\tchr 你画的图有些不准确，不过就定性而言，它确实表示了由两个等量异号电荷建立的场的电场线。你想象不出当其中一个电荷增大时，这幅图会怎样变化吗？

\stda 这种事我们从来没有做过。

\tchr 那就让我们来用场强矢量相加的法则。先从电荷相等这个熟悉的情形开始（\figr{fig-87}~\drkgry{(a)}）。选取 $A$、$B$、$C$ 三点，对每一点作一对场强矢量：$\vec{E}_{1}$ 和 $\vec{E}_{2}$（$\vec{E}_{1}$ 属于电荷 $+q_{1}$ 的场，$\vec{E}_{2}$ 属于电荷 $-q_{2}$ 的场）。然后对每一点把矢量 $\vec{E}_{1}$ 与 $\vec{E}_{2}$ 相加，得到合矢量 $\vec{E}_{A}$、$\vec{E}_{B}$ 和 $\vec{E}_{C}$。这些矢量应当与场的电场线在相应各点相切。这三个矢量指明了电场线的走向，如\figr{fig-88}~\drkgry{(a)} 所示。把这张图同你所画的\figr{fig-86} 比一比，注意你在电荷 $-q$ 左侧和电荷 $+q$ 右侧电场线走向上的不精确之处。现在设电荷 $+q_{1}$ 的大小加倍、电荷 $-q_{2}$ 减半（\figr{fig-87}~\drkgry{(b)}）。我们照旧取 $A$、$B$、$C$ 三点，先作出这些点上的场强矢量，再求出合矢量 $\vec{E}_{A}$、$\vec{E}_{B}$ 和 $\vec{E}_{C}$。与这些矢量对应的电场线图形如\figr{fig-88}~\drkgry{(b)} 所示。
\pnote{底本\figr{fig-87} 与\figr{fig-88} 的图题与\figr{fig-86} 的完全相同（且语法残缺），显系录入时误抄而未随图面改写：\figr{fig-87} 画的是场强矢量叠加的作图过程，\figr{fig-88} 画的是所得的电场线图形。中文版按图面内容分别拟写两图图题，待考原书。同段底本正文对图 86、87、88 的引用误写作公式引用（eq-86 等），恰好撞上第 15 章的同名公式标签，输出碰巧“看似正常”，中文版径改为图引用；又因原文件中\figr{fig-88} 的图块排在\figr{fig-87} 之前，两图的显示编号与文件名互换（与英文印本一致），不影响按图号的引用。本章另有若干标点级排印小误（must a be equal、takes the form , then、re- turn、Crosssections 等），径改不注。}

\begin{figure}[!h]%
\centering
\includegraphics[width=0.4\textwidth]{figs/ch-23/fig-088a}
\caption[与图 88 三种情形对应的电场线图形。]{与\figr{fig-87} 三种情形（a）（b）（c）对应的电场线图形。\label{fig-88}}
\end{figure}
\begin{figure*}%[-5cm]%
\centering
\includegraphics[width=\textwidth]{figs/ch-23/fig-087a}
\caption[用场强矢量叠加法确定电场线。]{用场强矢量叠加法确定两个异号点电荷的场的电场线：对 $A$、$B$、$C$ 三点分别作矢量 $\vec{E}_{1}$、$\vec{E}_{2}$ 并相加求合矢量。（a）两电荷等量；（b）$q_{1}$ 加倍、$q_{2}$ 减半；（c）再次加倍、减半。\label{fig-87}}
\end{figure*}
最后，设 $q_{1}$ 再次加倍、$q_{2}$ 再次减半（\figr{fig-87}~\drkgry{(c)}），对 $A$、$B$、$C$ 三点作出合矢量 $\vec{E}_{A}$、$\vec{E}_{B}$ 和 $\vec{E}_{C}$。相应的电场线图形如\figr{fig-88}~\drkgry{(c)} 所示。

正如你所看到的，随着电荷 $+q_{1}$ 相对大小的增大，它的影响变得更强：电荷 $+q_{1}$ 的场开始压制电荷 $-q_{2}$ 的场。

\stda 现在我明白了，怎样为多个电荷组成的系统所建立的场作电场线图。

\tchr 我们继续讨论静电场。这种场有一个重要性质，使它与引力场归为同一类，即：场力沿任何闭合路径所做的功等于零。换句话说，如果电荷在场中移动后返回出发点，那么在这一运动过程中场力做的功等于零。在路径的某些路段上这个功是正的，在另一些路段上是负的，但所做的功的总和等于零。由静电场的这一性质可以得出一些很有意思的推论。你能说出这些推论吗？

\stdb 说不上来，我想不出来。

\tchr 我来帮你。你大概已经注意到，静电场的电场线永远不会自行闭合。它们起止于电荷（始于正电荷、终于负电荷），或者伸向无穷远（或来自无穷远）。你能把这一情况与上面提到的静电场的性质联系起来吗？

\stdb 现在我似乎明白了。如果静电场中有一条电场线自行闭合，那么沿着它移动就能回到出发点。电荷沿电场线移动时，场做功的符号显然不变，因而这个功不可能等于零。可是另一方面，沿任何闭合路径所做的功都必须等于零。所以，静电场的电场线不可能自行闭合。
\pnote{底本此处仅有换行、漏排说话人标记（英文印本同样无标签）。按上下文，这是学生在回答教师刚提出的问题，当为之前答“说不上来”的学生乙，中文版补回，待考原书。}
\pnote{这段论证表明静电场是保守场（做功与路径无关的场）；大学普通物理把它写成环路定理 $\oint \vec{E} \cdot d\vec{l} = 0$。重力场同属保守场——重力做功与路径无关，由此引入重力势能（人教版必修第二册“机械能守恒定律”一章）；静电场中则相应引入电势能和电势（人教版必修第三册第十章“静电场中的能量”）。}

\tchr 完全正确。由静电场的上述性质还可以得出一个推论：把电荷从场中一点移到另一点所做的功，与所经过的路径无关。

我们可以把电荷从点 $a$ 移到点 $b$，譬如沿不同的路径 1 和 2（\figr{fig-89}）。用 $A_{1}$ 表示场力沿路径 1 移动电荷所做的功，用 $A_{2}$ 表示沿路径 2 所做的功。现在走一个完整的回路：沿路径 1 从点 $a$ 到点 $b$，再沿路径 2 从点 $b$ 返回点 $a$。沿路径 2 返回时，所做的功将是 $-A_{2}$。整个回路中做的总功为 $A_{1}+(-A_{2})=A_{1}-A_{2}$。既然沿任何自行闭合的路径所做的功都等于零，那么 $A_{1}=A_{2}$。移动电荷所做的功与所选路径无关、只与起点和终点有关——这一事实使这个量可以用来作为场的特征量（因为它只取决于场中所选的点！）。于是就引入了静电场的另一个特征量——电势（potential）。与场强不同，电势是标量，因为它是用功来表达的。

\begin{figure}[!h]%
\centering
\includegraphics[width=0.4\textwidth]{figs/ch-23/fig-089a}
\caption[静电场中做的功。]{静电场中做的功。\label{fig-89}}
\end{figure}
\stdb 我们在中学里学过：场的电势这个概念没有物理意义，只有场中任意两点的电势差（potential difference）才有物理意义。

\tchr 你说得完全对。严格说来，前面的讨论使我们能够引入的恰恰是电势差；场中 $a$、$b$ 两点之间的电势差（记作 $\left( \varphi_{a}-\varphi_{b} \right)$）定义为场力把电荷 $q_{0}$ 从点 $a$ 移到点 $b$ 所做的功与电荷 $q_{0}$ 之比，即
\begin{equation}%
\varphi_{a}-\varphi_{b}  = \frac{A_{a \to b}}{q_{0}}
\label{eq-132}
\end{equation}
而如果我们认为在无穷远处场不存在（即 $\varphi_{\infty}=0$），式~\eqref{eq-132} 就化为
\begin{equation}%
\varphi_{a} = \frac{A_{a \to \infty}}{q_{0}}
\label{eq-133}
\end{equation}
这样，场中给定点的电势，就可以用场力把单位正电荷从该点移到无穷远所做的功来确定。如果认为功不是场做的，而是反抗场力做的，那么给定点的电势就是把单位正电荷从无穷远移到该点所需做的功。自然，这样的定义排除了从实验上测量场中给定点电势的可能，因为在实验中我们无法退到无穷远去。正是由于这个原因，人们才说场中两点的电势差具有物理意义，而某点电势本身没有。可以说，给定点的电势只确定到相差一个任意常数的程度。通常就取无穷远处的电势作为这个常数，电势从这个值开始计量。为方便起见，假定无穷远处的电势等于零。
\pnote{原书“电势本身没有物理意义、电势差才有”的提法，对应人教版必修第三册第十章的零电势点选取：电势的数值取决于零参考点的选取（通常取无穷远处或大地为零电势点），具有确定物理意义的是电势差 $U_{AB}=\varphi_{A}-\varphi_{B}$，即本节式~\eqref{eq-132}。式~\eqref{eq-134} 的 $\varphi=q_{1}/r$ 仍是高斯制形式，SI 中点电荷的电势为 $\varphi=kQ/r$。}

在这些假定之下，点电荷 $q_{1}$ 建立的场在离电荷 $r$ 处一点的电势等于
\begin{equation}%
\varphi\,(r) = \frac{q_{1}}{r}
\label{eq-134}
\end{equation}
求几个点电荷建立的场在某点 $\vec{r}$ 处的电势，你们应该不会有什么困难。

\stdb 我们把每个电荷单独在点 $\vec{r}$ 产生的电势值分别记作 $\varphi_{1}(\vec{r})$、$\varphi_{2}(\vec{r})$ 等等。总电势 $\varphi (\vec{r})$ 显然等于各个电荷产生的电势的代数和。于是
\begin{equation}
\varphi (\vec{r}) = \varphi_{1}(\vec{r}) + \varphi_{2}(\vec{r}) + \ldots
\label{eq-135}
\end{equation}
求和时，正电荷产生的电势取正号，负电荷产生的电势取负号。

\tchr 完全正确。现在我们来考察等势面（equipotential surface）的概念。场中具有相同电势的点的轨迹称为等势面（或恒定电势面）。场中每一点都有一条电场线和一个等势面通过。它们彼此是怎样取向的？

\stdb 我知道，在每一点上电场线与等势面互相垂直。

\tchr 你能证明这一点吗？

\stdb 恐怕证不出来。

\begin{figure}[!h]%
\centering
\includegraphics[width=0.4\textwidth]{figs/ch-23/fig-090a}
\caption[证明电场线与等势面互相垂直。]{证明电场线与等势面互相垂直。\label{fig-90}}
\end{figure}
\tchr 这个证明并不难。设电场线 $aa_{1}$ 和等势面 $S$（\figr{fig-90}）都通过某一点 $a$。点 $a$ 处的场强由矢量 $\vec{E}_{a}$ 描述。现在把电荷 $q_{0}$ 从点 $a$ 移到等势面 $S$ 上离点 $a$ 很近、距离为 $\Delta l$ 的某一点 $b$。这一移动所做的功由下式表示
\begin{equation}%
A = F_{e}\Delta l \cos \alpha = E_{a}\, q_{0}\, \Delta l \cos \alpha
\label{eq-136}
\end{equation}
其中 $\alpha$ 是矢量 $\vec{E}_{a}$ 与移动方向之间的夹角。同一个功也可以表示为场在点 $a$ 和点 $b$ 的电势差。于是可以写出另一个关系式：
\begin{equation}%
A=q_{0} \left( \varphi_{a}-\varphi_{b} \right)
\label{eq-137}
\end{equation}
既然点 $a$ 与点 $b$ 位于同一个等势面上，就有 $\varphi_{a}=\varphi_{b}$。这就是说，根据式~\eqref{eq-137}，功 $A$ 应当等于零。把这个结果代入式~\eqref{eq-136}，得
\begin{equation}%
E_{a}\, q_{0} \Delta l \cos \alpha=0
\label{eq-138}
\end{equation}
式~\eqref{eq-138} 左边各因子中，只有 $\cos \alpha$ 可以等于零。于是我们得出结论：$\alpha=\ang{90}$。我想你们是清楚的：只要移动 $ab$ 的各种方向不越出等势面 $S$ 的范围，这一结果对任何移动方向都成立。表面的弯曲不妨碍我们的论证，因为位移 $\Delta l$ 非常小。除了使用电场线之外，人们还用等势面的截面来直观地描绘静电场。利用电场线与等势面互相垂直这一事实，可以从已知的电场线族画出等势面截面族，或者反过来。
\pnote{“电场线与等势面互相垂直”在人教版必修第三册第十章中作为结论直接给出（电场线由电势高的等势面指向电势低的等势面），本节补出了证明。证明中用的 $A=q_{0}\left( \varphi_{a}-\varphi_{b} \right)$ 即教材中静电力做功与电势差的关系 $W_{AB}=qU_{AB}$。导体表面是等势面、因而电场线与导体表面垂直的推论，见\chap{ch-24}。}

（问学生甲）你来试着为\figr{fig-88}~\drkgry{(a)} 所示的情形画一画等势面的截面，好吗？为了不与电场线混淆，请把等势面的截面画成虚线。

\stda 我画虚线时，让它们总是与电场线成直角相交。这就是我画的图（\figr{fig-91}）。

\begin{figure}%[-5cm]%
\centering
\includegraphics[width=0.5\textwidth]{figs/ch-23/fig-091a}
\caption[等势面的截面与电场线。]{等势面的截面与电场线。\label{fig-91}}
\end{figure}
\tchr 你画得对。

\end{dialogue}



\cleardoublepage

% NEWTERMS: electrostatic field | 静电场 | 章题与全章高频；底本该词于 ch-22（待译）先出现，首现括注归属以成书为准，届时可撤一处
% NEWTERMS: intensity of the electric field | 电场强度 | 简称场强；符号 $E$；人教版必修第三册第九章“电场 电场强度”；首现括注；高斯制形式见评注
% NEWTERMS: lines of force | 电场线 | 直译“力线”；人教版必修第三册通名“电场线”；首现括注
% NEWTERMS: point charge | 点电荷 | 人教版必修第三册第九章；首现括注
% NEWTERMS: Coulomb force | 库仑力 | 库仑定律见人教版必修第三册第九章“库仑定律”一节；原书高斯制形式见评注
% NEWTERMS: potential | 电势 | 原书符号 $\varphi$；人教版必修第三册第十章“电势能和电势”；首现括注
% NEWTERMS: potential difference | 电势差 | 人教版记 $U_{AB}$，对应原书 $\varphi_{a}-\varphi_{b}$；首现括注
% NEWTERMS: equipotential surface | 等势面 | 原书别名 surface of constant potential（恒定电势面）；首现括注
% NEWTERMS: cross-sections of equipotential surfaces | 等势面的截面 | 纸面内的截线即等势线；图 91 语境
% NEWTERMS: resultant intensity | 合场强 | 场强矢量叠加语境
